\documentclass[11pt]{article}
\usepackage{amsmath,amssymb,amsthm}
\usepackage[margin=1.1in]{geometry}
\usepackage{booktabs}
\usepackage{hyperref}

\newtheorem{proposition}{Proposition}
\newtheorem{remark}{Remark}

\title{Measuring Success:\\
Confusion Matrices, Curves, and the Cost of a Threshold}
\author{Yuval Lavie}
\date{\today}

\begin{document}
\maketitle

\begin{abstract}
Training minimizes a loss; \emph{evaluation} answers a different
question --- did the classifier do its job? --- and the answer depends
on which job. This page organizes the metric zoo around three ideas:
every metric is a function of the four confusion-matrix counts; the
threshold is a free parameter, so honest evaluation sweeps it (ROC,
precision--recall); and calibrated probabilities turn misclassification
\emph{costs} into the optimal threshold in closed form. Every claim is
verified by \texttt{evaluation\_demo.py} on held-out sets of $2$--$4
\times 10^5$ points, and the animation makes the central fact visible:
\emph{every point of the ROC and PR curves is a threshold}.
\end{abstract}

\section{The four counts}

A binary classifier in this repository is a score $s(x)$ --- for the
GLM family, a calibrated probability $\hat p(x)$ --- plus a threshold
$\tau$; it predicts $+$ iff $s(x) > \tau$. Against the truth, four
counts:
\begin{center}
\begin{tabular}{lcc}
\toprule
 & predicted $+$ & predicted $-$\\
\midrule
truly $+$ & TP & FN \ (misses)\\
truly $-$ & FP \ (false alarms) & TN\\
\bottomrule
\end{tabular}
\end{center}
Everything below is arithmetic on these four numbers:
\begin{equation}
  \mathrm{acc} = \frac{TP+TN}{n}, \quad
  \mathrm{prec} = \frac{TP}{TP+FP}, \quad
  \mathrm{rec} = \mathrm{TPR} = \frac{TP}{TP+FN}, \quad
  \mathrm{FPR} = \frac{FP}{FP+TN},
  \label{eq:defs}
\end{equation}
plus specificity $= 1 - \mathrm{FPR}$ and the harmonic compromise
\begin{equation}
  F_1 = \frac{2\,\mathrm{prec}\cdot\mathrm{rec}}
             {\mathrm{prec}+\mathrm{rec}},
  \qquad
  F_\beta = (1+\beta^2)\,
  \frac{\mathrm{prec}\cdot\mathrm{rec}}
       {\beta^2\,\mathrm{prec}+\mathrm{rec}}
  \label{eq:f1}
\end{equation}
($\beta > 1$ weights recall; harmonic, so one bad component drags the
mean down --- you cannot buy $F_1$ with precision alone). Each metric
answers a \emph{different question}: accuracy ``how often right'';
precision ``when it cries $+$, believe it?''; recall ``of the real
$+$, how many found?''.

\paragraph{Accuracy lies under imbalance.} Measured, base rate
$2.67\%$: the classifier that \emph{never says positive} scores
accuracy $0.9732$ with recall $0.000$; the fitted logistic scores
$0.9762$ with recall $0.212$. Three thousandths of accuracy separate
``useless'' from ``useful'' --- accuracy answers the wrong question
whenever one class is rare and expensive to miss. This is why the rest
of this page exists.

\section{The threshold is a free parameter}

The score is the model; $\tau$ is a \emph{decision}. Sweeping $\tau$
from $1$ to $0$ trades misses for false alarms and traces two curves:
\begin{itemize}
  \item \textbf{ROC}: $\bigl(\mathrm{FPR}(\tau),
    \mathrm{TPR}(\tau)\bigr)$ --- both coordinates are per-class rates.
  \item \textbf{Precision--recall}:
    $\bigl(\mathrm{rec}(\tau), \mathrm{prec}(\tau)\bigr)$ --- precision
    mixes the classes.
\end{itemize}
Computation note (a bug the asserts caught): the exact curve comes from
\emph{sorting by score} --- each prefix of the sorted list is one
threshold --- not from a grid of $\tau$ values, which misses the
saturated tails where sigmoid scores pile up near $0$ and $1$.

\begin{proposition}[What the area under ROC is]
\begin{equation}
  \boxed{\;\mathrm{AUC} \;=\; \Pr\bigl(s(X^+) > s(X^-)\bigr)
  \;+\; \tfrac12 \Pr\bigl(s(X^+) = s(X^-)\bigr)\;}
  \label{eq:auc}
\end{equation}
for independent draws $X^+$ from the positive class and $X^-$ from the
negative class.
\end{proposition}

\begin{proof}[Proof sketch]
$\mathrm{AUC} = \int_0^1 \mathrm{TPR}\, d\mathrm{FPR}$. Substituting
$\mathrm{FPR}(\tau) = \Pr(s(X^-) > \tau)$ as the integration variable,
$d\mathrm{FPR} = -f_{s^-}(\tau)\,d\tau$ and $\mathrm{TPR}(\tau) =
\Pr(s(X^+) > \tau)$, so the integral is
$\int \Pr(s(X^+) > \tau)\, f_{s^-}(\tau)\, d\tau
= \Pr(s(X^+) > s(X^-))$ by conditioning on $s(X^-) = \tau$. (Ties
contribute the diagonal half-steps.)
\end{proof}

Measured: exact-curve area $0.8957$ vs.\ $5\times10^5$ Monte Carlo
score pairs $0.8954$. AUC is a pure \emph{ranking} quality --- blind
to calibration, blind to the threshold, blind to the class ratio.

\paragraph{ROC is immune to imbalance; PR is not.} Deleting $95\%$ of
the positives (score distributions untouched) moves AUC by $0.003$
($0.8957 \to 0.8984$) but collapses average precision from $0.8585$ to
$0.4537$: precision divides by \emph{predicted positives}, which
rebases on the class ratio. Under heavy imbalance, the PR curve is
where the pain shows --- read it.

\section{Calibration turns costs into thresholds}

A score is \emph{calibrated} if $\Pr(Y{=}1 \mid s(X) = p) = p$ ---
the score means what it says. The well-specified logistic MLE is
calibrated by construction (log-loss's population minimizer is the
true conditional probability; Theory/losses): measured reliability gap
$\le 0.011$ per decile, Brier score $0.129$.

Calibration is what makes decisions computable. Let a false alarm cost
$c_{FP}$ and a miss cost $c_{FN}$. Predicting $+$ at a point with
probability $p$ has expected cost $(1-p)\,c_{FP}$; predicting $-$
costs $p\,c_{FN}$. Predict $+$ exactly when the first is smaller:
\begin{equation}
  \boxed{\;\tau^\ast \;=\; \frac{c_{FP}}{c_{FP} + c_{FN}}\;}
  \label{eq:tau}
\end{equation}
--- the optimal threshold is a \emph{ratio of costs}, no sweep
required. Measured with $c_{FN} = 5\,c_{FP}$: theory $\tau^\ast =
0.167$, empirical cost minimum at $0.161$. Two corollaries worth
keeping: symmetric costs give $\tau^\ast = 1/2$ (the Bayes classifier
of Algorithms/classification/binary/logistic/); and none of this works if the scores are
uncalibrated --- an SVM's margins (Theory/losses: hinge remembers only
the sign) cannot be thresholded by \eqref{eq:tau} without a
calibration step first.

\paragraph{$F_1$'s hidden threshold.} $F_1$ is also maximized at a
data-dependent threshold, not $0.5$: measured $0.231$ on the
imbalanced problem. If a metric picks your threshold, know which
question that metric answers.

\section{Choosing, in one table}

\begin{center}
\begin{tabular}{lll}
\toprule
metric & the question it answers & fails when\\
\midrule
accuracy & how often right & imbalance (base-rate dominance)\\
precision & believe its $+$? & ignores misses entirely\\
recall / TPR & how many $+$ found & ignores false alarms\\
$F_1$/$F_\beta$ & prec--rec compromise & costs aren't sym./explicit\\
ROC--AUC & ranking quality & imbalance hides FP pain; no $\tau$\\
PR--AP & rare-class usefulness & incomparable across base rates\\
log-loss / Brier & probability quality & needs calibrated scores\\
expected cost & the actual decision & needs costs + calibration\\
\bottomrule
\end{tabular}
\end{center}

\section{Experiment}

All numbers above from \texttt{evaluation\_demo.py}: logistic MLEs
fitted by IRLS on a balanced ($56.5\%$) and an imbalanced ($2.67\%$)
version of the standard 2D problem, evaluated on $2$--$4\times10^5$
held-out points; every quantitative claim is an assert. The static
figure shows the reliability diagram, the cost-vs-threshold curve with
$\tau^\ast$ marked, and the PR collapse under imbalance. The animation
slides $\tau$ from $0$ to $1$: the decision line shifts in parallel
across the probability field (same $w$, different level set), the
confusion report updates live, and two dots travel the ROC and PR
curves --- every point of those curves \emph{is} a threshold.

One sentence to keep: \emph{the model earns the score, but the
threshold spends it --- and only calibrated probabilities let you
spend optimally.}

\end{document}
